% ============================================================
%  Section 2 — Related Work  (~1 page)
% ============================================================
\section{Related Work}

\subsection{Framework-Specific Benchmarking}
A substantial body of work characterises the performance of a single
quantum SDK in isolation.
Studies of Qiskit's transpiler optimisation levels show that higher
optimisation passes reduce CNOT counts by 10--40\% on typical circuits,
but these evaluations run exclusively on Qiskit's Aer simulator and
cannot be compared directly to equivalent Cirq or PennyLane measurements.
Similarly, PennyLane benchmarks in the variational literature report
gradient-evaluation throughput on its \texttt{default.qubit} backend, a
metric with no direct counterpart in Qiskit or Cirq.
Framework-specific benchmarks are valuable for intra-framework
optimisation studies; they do not establish cross-framework rankings.

\subsection{Cross-Framework Comparisons}
Multi-framework studies exist but carry the simulator confound described
above.
The most common experimental design pairs each framework with its native
backend (Qiskit + Aer, Cirq + cirq-google simulator,
PennyLane + default.qubit) and reports fidelity or runtime figures side
by side.
An observed difference in simulation time between Qiskit and Cirq in
such a study conflates transpiler latency, compiled circuit size, and the
throughput of two distinct simulators---three variables that cannot be
disentangled post hoc.
A second limitation of existing cross-framework studies is scope: most
restrict comparison to one or two algorithms (typically Bell state and
Quantum Fourier Transform) without systematically varying qubit count,
which precludes scaling analysis.
To our knowledge, this paper presents the first study that routes all
circuits from all frameworks through a single common backend before
comparison.

\subsection{Quantum Volume and Holistic Benchmarks}
Cross et~al.\ introduced Quantum Volume (QV) as a single-number
estimator of a device's effective circuit-processing
capacity~\cite{cross2019validating}.
QV integrates gate fidelity, connectivity, and compilation quality into
one figure, making it attractive as a holistic comparison metric.
QPack-style composite scores extend this idea to multi-dimensional
profiles that separate compilation speed, scalability, accuracy, and
capacity into independently interpretable sub-scores.
We apply QV in a structure-implied, simulator-agnostic mode using a
depolarizing noise model~\cite{nielsen2010quantum}, which differs from
IBM's original hardware-measured QV but allows fair cross-framework
comparison on a common backend.
